\section{Behavior Under Algorithmic Redistricting}
\label{sec:behavior}

\subsection{Why Bisect Reduces MM}

\textsc{Bisect} reduces MM relative to enacted maps through the same compactness
mechanism that reduces EG, but the mechanism operates differently at the level
of the vote-share distribution.

Under enacted maps drawn with partisan intent, packing concentrates Democratic
voters in a small number of districts with Democratic vote shares of 75--85\%.
These packed districts pull the mean Democratic vote share substantially above
the median, because the median is determined by the cluster of competitive or
Republican-leaning districts in the middle of the distribution, not by the
packed outliers. The result is large positive MM.

Under \textsc{bisect}, compactness constraints prevent the geographic reaching
required to concentrate Democratic voters from dispersed locations into a small
number of super-majority districts. Compact districts centered on urban cores
will have high Democratic vote shares---the core concentration is a geographic
fact---but the absence of deliberate additional packing means fewer districts
at the extreme high end of the distribution. The distribution of Democratic
vote shares is less right-skewed, and the mean-median gap is smaller.

The residual MM in \textsc{bisect} maps ($\approx 0.02$--$0.03$) reflects the
minimum right-skewness that results from compact districts centered on North
Carolina's Triangle and Charlotte urban cores. These cores have Democratic vote
shares of 70--80\%, and any compact district that includes them will have
Democratic vote shares near that range. Compactness prevents drawing additional
territory to push adjacent districts from 48\% to 55\% Democratic---the cracking
operation that would add more moderate-margin districts to the left of the
median---but it cannot eliminate the right-tail outliers created by the urban cores
themselves.

\subsection{The Geographic Sorting Baseline}

The \textsc{bisect} MM value of $\approx 0.02$--$0.03$ for NC represents the
geographic sorting baseline: the MM that any compact, population-equal redistricting
plan must produce for North Carolina given the residential distribution of its
voters. This baseline is empirically determined by running \textsc{bisect} on
the state's census tract geometry with VEST election returns.

The baseline is state-specific: Texas ($k=38$) has a slightly higher baseline
($\approx 0.03$) because its urban cores (Houston, Dallas, Austin, San Antonio)
are larger and more densely Democratic relative to the state's overall partisan
balance. Wisconsin ($k=8$) has a similar baseline ($\approx 0.02$) because
Milwaukee and Madison are dense Democratic centers relative to the state's
competitive statewide balance.

\subsection{The Mapmaker Contribution}

The mapmaker contribution to MM is the difference between enacted MM and
\textsc{bisect} MM, holding geography constant:
\[
\Delta\text{MM} = \text{MM}_{\text{enacted}} - \text{MM}_{\text{bisect}}.
\]
For NC: $\Delta\text{MM} \approx 0.07 - 0.02 = 0.05$. This gap represents
the additional vote-share skewness that mapmakers introduced by deliberately
drawing boundaries to concentrate Democratic voters in fewer, higher-margin
districts than compactness would produce naturally. In state constitutional
terms, $\Delta\text{MM}$ is the portion of partisan advantage that cannot
be explained by geography and therefore must be explained by deliberate
redistricting choices.

This decomposition is the key legal contribution of comparing \textsc{bisect}
to enacted maps. The geographic component ($\approx 0.02$) is irreducible;
courts applying NC's free elections clause should focus on the mapmaker
component ($\approx 0.05$) as the portion attributable to deliberate manipulation.
